\documentclass[10pt]{article}

% ---------- Document Identity (update these only) ----------
\newcommand{\DocStage}{}
\newcommand{\DocTitle}{Your Road to Tier One SOF}
\newcommand{\ProgramName}{182-Week Civilian-to-Selection Readiness Pipeline}
\newcommand{\ProgramSubtitle}{}

% ---------- Page ----------
\usepackage[legalpaper,left=0.5in,right=0.5in,top=0.6in,bottom=0.45in]{geometry}

% ---------- Typography ----------
\usepackage[T1]{fontenc}
\usepackage[utf8]{inputenc}
\usepackage{libertinus}
\usepackage{microtype}
\usepackage{parskip}
\usepackage{ragged2e}
\usepackage{textcomp}
\AtBeginDocument{\color{black}}
\AtBeginDocument{\pagecolor{white}}
\AtBeginDocument{\justifying}

% ---------- Landscape pages ----------
\usepackage{pdflscape}

% ---------- Color ----------
\usepackage[table]{xcolor}
\definecolor{StageBlue}{HTML}{1F4E79}
\definecolor{StageBlueDark}{HTML}{163A5A}
\definecolor{LightGray}{HTML}{F2F4F7}
\definecolor{MidGray}{HTML}{D9DEE5}
\definecolor{RuleGray}{HTML}{C7CED8}
\definecolor{HeaderInk}{HTML}{000000}
\definecolor{Ink}{HTML}{000000}

% ---------- Utility ----------
\usepackage{enumitem}
\sloppy
\setlength{\emergencystretch}{2em}

% ---------- Boxes / UI ----------
\usepackage[most]{tcolorbox}

\tcbset{
  charterTitle/.style={
    enhanced,
    colback=StageBlue!4,
    colframe=StageBlue,
    boxrule=1.25pt,
    arc=3pt,
    left=16pt,right=16pt,top=14pt,bottom=14pt,
    borderline north={3.2pt}{0pt}{StageBlue},
    borderline west={4pt}{0pt}{StageBlue},
    borderline south={1.25pt}{0pt}{StageBlue},
    drop shadow={black!25!white},
  },
  charterRule/.style={
    enhanced,
    colback=LightGray,
    colframe=RuleGray,
    boxrule=0.85pt,
    arc=3pt,
    left=12pt,right=12pt,top=10pt,bottom=10pt,
    borderline west={3pt}{0pt}{StageBlue},
  },
  charterBadge/.style={
    enhanced,
    colback=StageBlue,
    coltext=white,
    colframe=StageBlueDark,
    boxrule=0.6pt,
    arc=2pt,
    left=6pt,right=6pt,top=3pt,bottom=3pt,
  }
}
\newtcbox{\Badge}{nobeforeafter,charterBadge}

% ---------- Lists ----------
\setlist[itemize]{leftmargin=1.5em, itemsep=0.25em, topsep=0.25em}
\setlist[enumerate]{leftmargin=1.8em, itemsep=0.25em, topsep=0.25em}

% ---------- TOC ----------
\usepackage{tocloft}
\setcounter{tocdepth}{3}
\setlength{\cftbeforesecskip}{3pt}
\setlength{\cftbeforesubsecskip}{2pt}
\setlength{\cftbeforesubsubsecskip}{0pt}
\renewcommand{\cftsecfont}{\bfseries}
\renewcommand{\cftsecpagefont}{\bfseries}
\renewcommand{\cftsubsecfont}{\normalfont}
\renewcommand{\cftsubsecpagefont}{\normalfont}
\renewcommand{\cftsubsubsecfont}{\normalfont}
\renewcommand{\cftsubsubsecpagefont}{\normalfont}
\setlength{\cftsecindent}{0pt}
\setlength{\cftsubsecindent}{1.25em}
\setlength{\cftsubsubsecindent}{2.8em}
\setlength{\cftsecnumwidth}{2.1em}
\setlength{\cftsubsecnumwidth}{2.8em}
\setlength{\cftsubsubsecnumwidth}{3.6em}

% Remove the built-in TOC heading so only the custom heading is shown
\makeatletter
\renewcommand{\tableofcontents}{\@starttoc{toc}}
\makeatother

% ---------- Header/Footer: page number only ----------
\usepackage{fancyhdr}
\pagestyle{fancy}
\fancyhf{}
\renewcommand{\headrulewidth}{0pt}
\renewcommand{\footrulewidth}{0pt}
\fancyfoot[C]{\thepage}

% ---------- Section formatting ----------
\usepackage{titlesec}

\titleformat{\section}
  {\Large\bfseries\color{Ink}}
  {\thesection}{0.6em}{}
  [\vspace{0.25em}\textcolor{StageBlue}{\rule{\linewidth}{0.8pt}}]

\titleformat{\subsection}
  {\large\bfseries\color{Ink}}
  {\thesubsection}{0.6em}{}

\titleformat{\subsubsection}
  {\normalsize\bfseries\color{Ink}}
  {\thesubsubsection}{0.6em}{}

\titlespacing*{\section}{0pt}{0.95em}{0.35em}
\titlespacing*{\subsection}{0pt}{0.65em}{0.25em}
\titlespacing*{\subsubsection}{0pt}{0.55em}{0.2em}

% ---------- Table engine ----------
\usepackage{tabularray}
\UseTblrLibrary{booktabs}

% ---------- tabularray: GLOBAL table treatment ----------
\SetTblrInner{
  cells = {fg=black, bg=white, valign=t},
}

% ---------- Header helpers ----------
\newcommand{\Hdr}[1]{\shortstack[c]{\strut #1\strut}}

% ---------- Lines ----------
\newcommand{\TitleLine}{\textcolor{StageBlue}{\rule{0.98\linewidth}{0.95pt}}}
\newcommand{\SubTitleLine}{\textcolor{RuleGray}{\rule{0.98\linewidth}{0.65pt}}}

% ---------- Continued on next page ----------
\newcommand{\ContinuedOnNextPageInline}{%
  \par
  \vfill
  \noindent\textcolor{RuleGray}{\rule{\linewidth}{1pt}}\vspace{-2pt}
  \noindent\hfill\textit{Continued on next page}\par\vspace{-10pt}
  \noindent\textcolor{RuleGray}{\rule{\linewidth}{1pt}}
  \clearpage
}

% ---------- PDF hyperlinks ----------
\usepackage[hidelinks]{hyperref}
\hypersetup{
  colorlinks=false,
  pdfborder={0 0 0},
  linktoc=all,
  pdftitle={\DocTitle},
  pdfauthor={\ProgramName}
}

% =========================================================
\begin{document}

\subsection*{Phase 11 Card Header}
\phantomsection
\addcontentsline{toc}{subsection}{Phase 11 Card Header}

\begin{itemize}
\item Phase \textbf{ID}: Phase 11
\item Phase Name: Pre-Selection Conditioning and Recovery
\item Duration weeks: 12
\item Version: v1.0
\item Stage Alignment: Stage 5
\item Governing Status: Draft — Non-Deployable
\end{itemize}

\subsection*{1. Phase Mission}
\phantomsection

Phase 11 consolidates pre-selection readiness by prioritizing verification-quality evidence under conservative recovery control, using protected windows to stage high-consequence tests without stacking, and enforcing reslot discipline when validity or safety is compromised.

\subsection*{2. Phase Purpose}
\phantomsection

\begin{itemize}
\item Establish repeatable, gate-grade performance evidence that is stable under controlled fatigue and recovery posture.
\item Close replication gaps for gate-critical metrics without overtesting, using protected windows and staged batteries.
\item Prepare Phase 12 by proving that high-consequence events (especially long \textbf{RUN} and long \textbf{RUCK}) can be produced as \textbf{VALID} evidence without durability collapse.
\end{itemize}

\subsection*{3. Entry Criteria}
\phantomsection

\begin{itemize}
\item Entry requires admissible evidence that Phase 10 exit posture is satisfied under Stage 1.F admissibility rules (\textbf{VALID} tests, complete logs, no unresolved stop-rule events).
\item Default posture if evidence is incomplete or invalid: \textbf{HOLD} and reslot per Stage 3.E; do not proceed into Phase 11 as “assumed passed.”
\end{itemize}

\subsection*{4. Operating Constraints}
\phantomsection

\begin{itemize}
\item 6 sessions per week.
\item Required session elements per Stage 1.F in correct order:
  \begin{itemize}
  \item Safety self-check first
  \item Warm-up
  \item Mobility
  \item Main work
  \item Cardio
  \item Cooldown
  \item Recovery notes and any required post-session notes per Stage 1.F
  \end{itemize}
\item Cardio minimum standard: minimum 30 minutes continuous per session.
\item 10,000 steps per day global requirement.
\item Validity doctrine: admissible Valid evidence only for gates. \textbf{INVALID} tests provide no gate credit and must be reslotted per Stage 3.E.
\end{itemize}

\subsection*{5. Primary Adaptations and Physiological Focus}
\phantomsection

\begin{itemize}
\item High-specificity locomotion readiness with low variance outputs under controlled fatigue.
\item Durability stability under higher consequence domains (\textbf{RUN} and \textbf{RUCK}) without gait drift or foot breakdown.
\item Recovery sensitivity management: sleep, soreness, fatigue, and pain trends constrain testing attempts and trigger deload-only conversions when needed.
\item Calisthenics repeatability preserved under verification posture.
\item Evidence integrity: stable environments, stable tools, and complete logging to preserve admissibility and replication.
\end{itemize}

\subsection*{6. Primary Modalities}
\phantomsection

\begin{itemize}
\item Emphasized domains: \textbf{RUN}; \textbf{RUCK}; \textbf{CAL}; \textbf{DURABILITY}; \textbf{RECOVERY}.
\item Supporting domains (phase-appropriate, not substitute evidence): \textbf{SWIM}; \textbf{STR}; \textbf{BODYCOMP}; \textbf{AER}.
\item De-emphasized/prohibited:
  \begin{itemize}
  \item No invention of tests or thresholds.
  \item Underwater (\textbf{C14} / \textbf{MET-2B-025} and \textbf{MET-2B-026}) is conditional and governance-gated; if governance cannot be guaranteed, it is not scheduled and cannot be used for gate credit.
  \end{itemize}
\end{itemize}

\clearpage
\begin{landscape}

\subsection*{7. Weekly Structure}
\phantomsection

\begingroup
\scriptsize
\setlength{\tabcolsep}{2pt}
\renewcommand{\arraystretch}{1.12}

\begin{longtblr}{
  width=\linewidth,
  rowhead=1,
  colspec = {
    Q[l,wd=0.70in]
    X[l,2.10]
    X[l,1.55]
    X[l,2.95]
    X[l,2.90]
  },
  hlines,
  vlines,
  cells = {hsep=6pt, vsep=6pt, halign=l, valign=t},
  cell{1-Z}{1-5} = {fg=black},
  row{odd} = {bg=white},
  row{even} = {bg=LightGray},
  row{1} = {bg=StageBlue!15, fg=black, font=\bfseries, halign=l, valign=t},
}
\Hdr{Session \#} &
\Hdr{Primary Focus} &
\Hdr{Main Work Domains} &
\Hdr{Cardio Modality/Intent} &
\Hdr{Notes/Risk Controls} \\

1 &
\textbf{RUN} readiness exposure under controlled cost &
\textbf{RUN}; \textbf{DURABILITY} &
Modality: step-based locomotion or treadmill as needed; Minutes: 30–60; RPE plus Talk Test cue: RPE 3–5, full to short sentences; Continuity: continuous planned &
No novelty near test windows; protect gait integrity; document footwear and route/tool stability. \\

2 &
\textbf{CAL} repeatability and trunk endurance preservation &
\textbf{CAL}; \textbf{DURABILITY} &
Modality: low-impact steady; Minutes: 30–45; RPE plus Talk Test cue: RPE 3–4, full sentences; Continuity: continuous planned &
\textbf{CAL} gate artifacts required when testing; avoid overuse flare; prioritize technique and auditability. \\

3 &
\textbf{RUCK} tolerance maintenance under strict foot governance &
\textbf{RUCK}; \textbf{DURABILITY} &
Modality: low-impact steady; Minutes: 30–60; RPE plus Talk Test cue: RPE 3–5, full to short sentences; Continuity: continuous planned &
Foot hotspot escalation triggers downshift; do not drift load/configuration near gate windows. \\

4 &
Recovery-biased consolidation and mobility emphasis &
\textbf{RECOVERY}; \textbf{DURABILITY} &
Modality: lowest-risk continuous modality; Minutes: 30–45; RPE plus Talk Test cue: RPE 2–4, full sentences; Continuity: continuous planned &
Acts as fatigue-control day; used to protect upcoming protected-window validity. \\

5 &
Multi-domain maintenance without stacking high-cost exposures &
\textbf{RUN} or \textbf{RUCK} (supporting); \textbf{CAL} (supporting); \textbf{RECOVERY} &
Modality: steady; Minutes: 30–60; RPE plus Talk Test cue: RPE 3–5, controlled; Continuity: continuous planned &
Do not create hidden test days; avoid pairing maximal \textbf{RUN} and long \textbf{RUCK} within spacing intent. \\

6 &
Readiness confirmation and administrative integrity day &
\textbf{DURABILITY}; \textbf{RECOVERY}; \textbf{BODYCOMP} &
Modality: steady; Minutes: 30–45; RPE plus Talk Test cue: RPE 2–4, full sentences; Continuity: continuous planned &
Log completeness check; address tool/environment drift before protected windows. \\

\end{longtblr}

\endgroup

\subsection*{8. Progression Rules}
\phantomsection

\begin{itemize}
\item No stacking make-ups: missed/invalid tests are reslotted to the next protected window per Stage 3.E; do not compress.
\item One progression lever at a time: do not increase multiple high-cost demands in the same week-in-phase.
\item If durability or recovery trends deteriorate (sleep collapse, rising pain, hotspots), convert planned test attempts to deload-only or check-in-only and reslot maximal tests.
\item Protect decisive proof events:
  \begin{itemize}
  \item If the phase includes a long \textbf{RUCK} milestone, treat it as a dominant event and do not pair with maximal \textbf{RUN} inside the same 96-hour intent window.
  \end{itemize}
\item Substitute to preserve safety, not to escalate intensity. Any substitution that increases risk is prohibited.
\end{itemize}

\subsection*{9. Deload and Test Placement}
\phantomsection

Stage 3 integrated calendar posture for Phase 11 is taper-aware. Protected windows by week-in-phase:

\begin{itemize}
\item Week 1 baseline touchpoint (low risk; not a maximal gate by default)
\item Week 5 Deload+Test (mid-phase decision touchpoint)
\item Week 10 Deload+Test (late-phase verification / long-load milestone window by intent)
\item Week 11 Transition (stabilization and retest buffer only; no new major tests)
\item Week 12 Deload+Test (end-of-phase exit gate window)
\end{itemize}

\end{landscape}
\clearpage

\begin{landscape}

\subsection*{10. Test Battery}
\phantomsection

\begingroup
\scriptsize
\setlength{\tabcolsep}{2pt}
\renewcommand{\arraystretch}{1.12}

\begin{longtblr}{
  width=\linewidth,
  rowhead=1,
  colspec = {
    X[l,1.45]
    X[l,3.15]
    Q[l,wd=0.95in]
    X[l,3.25]
    Q[l,wd=1.35in]
  },
  hlines,
  vlines,
  cells = {hsep=6pt, vsep=6pt, halign=l, valign=t},
  cell{1-Z}{1-5} = {fg=black},
  row{odd} = {bg=white},
  row{even} = {bg=LightGray},
  row{1} = {bg=StageBlue!15, fg=black, font=\bfseries, halign=l, valign=t},
}
\Hdr{Test Timing} &
\Hdr{Test \textbf{ID}/Name} &
\Hdr{Domain} &
\Hdr{Validity Conditions} &
\Hdr{Pass Threshold Stage 2 Tier} \\

Week 1 baseline &
\textbf{C1} / \textbf{MET-2B-001} Bodyweight — Weekly Average &
\textbf{BODYCOMP} &
Standard conditions and complete logging per Stage 2; method stability &
\textbf{INTERMEDIATE} \\

Week 1 baseline &
\textbf{C7} / \textbf{MET-2B-013} Plank — Forearm &
\textbf{CAL} &
\textbf{VALID} test conditions; \textbf{CAL} artifact posture per Stage 2; no stop-rule violations &
\textbf{INTERMEDIATE} \\

Week 5 mid-phase gate &
\textbf{C10} / \textbf{MET-2B-016} 3-Mile Run Time Trial &
\textbf{RUN} &
\textbf{VALID} conditions; route/tool stability per Stage 2.E; warm-up logged; no stop-rule violations &
\textbf{INTERMEDIATE} \\

Week 5 mid-phase gate &
\textbf{C4} / \textbf{MET-2B-010} Push-Ups — 2-Minute Timed, Strict &
\textbf{CAL} &
\textbf{VALID} conditions; \textbf{CAL} artifact posture per Stage 2; standardized counting &
\textbf{INTERMEDIATE} \\

Week 5 mid-phase gate &
\textbf{C5} / \textbf{MET-2B-011} Sit-Ups — 2-Minute Timed, Standard &
\textbf{CAL} &
\textbf{VALID} conditions; \textbf{CAL} artifact posture per Stage 2; standardized counting &
\textbf{INTERMEDIATE} \\

Week 5 mid-phase gate &
\textbf{C6} / \textbf{MET-2B-012} Pull-Ups — Strict Dead-Hang &
\textbf{CAL} &
\textbf{VALID} conditions; \textbf{CAL} artifact posture per Stage 2; strict standard &
\textbf{INTERMEDIATE} \\

Week 10 milestone window &
\textbf{C12} / \textbf{MET-2B-021} 12-Mile Ruck Time Trial — 55 lb &
\textbf{RUCK} &
\textbf{VALID} conditions; dry load verification; water logged separately; route/tool stability per Stage 2.E; no stop-rule violations &
\textbf{SELECTION-READY} \\

Week 12 end gate &
\textbf{C11} / \textbf{MET-2B-017} 5-Mile Run Time Trial &
\textbf{RUN} &
\textbf{VALID} conditions; route/tool stability; warm-up logged; no stop-rule violations; avoid novelty near gate &
\textbf{SELECTION-READY} \\

Week 12 end gate &
\textbf{C4} / \textbf{MET-2B-010} Push-Ups — 2-Minute Timed, Strict &
\textbf{CAL} &
\textbf{VALID} conditions; \textbf{CAL} artifact posture per Stage 2 &
\textbf{SELECTION-READY} \\

Week 12 end gate &
\textbf{C5} / \textbf{MET-2B-011} Sit-Ups — 2-Minute Timed, Standard &
\textbf{CAL} &
\textbf{VALID} conditions; \textbf{CAL} artifact posture per Stage 2 &
\textbf{SELECTION-READY} \\

Week 12 end gate &
\textbf{C6} / \textbf{MET-2B-012} Pull-Ups — Strict Dead-Hang &
\textbf{CAL} &
\textbf{VALID} conditions; \textbf{CAL} artifact posture per Stage 2 &
\textbf{SELECTION-READY} \\

Conditional, not scheduled unless governance guaranteed &
\textbf{C14} / \textbf{MET-2B-025} Underwater Exposure Admissibility Flag; \textbf{C14} / \textbf{MET-2B-026} Underwater 25 m Time &
\textbf{SWIM} &
Only if governance prerequisites satisfied (certified lifeguard + dedicated safety spotter + facility permission); otherwise not scheduled and inadmissible &
\textbf{SELECTION-READY} \\

\end{longtblr}

\endgroup

\end{landscape}
\clearpage

\subsection*{11. Exit Standards}
\phantomsection

\begin{itemize}
\item Gate outcome at end of Phase 11 is based only on admissible evidence:
  \begin{itemize}
  \item Tests must be \textbf{VALID}; invalid provides no gate credit and triggers reslot posture per Stage 3.E.
  \item Replication posture for tier classification is governed by Stage 2 (2 of last 3 \textbf{VALID} within decision window).
  \end{itemize}
\item Intended tier posture for Phase 11 exit: \textbf{SELECTION-READY}.
\item \textbf{ADVANCE} requires that Phase 11 end-gate tests meet the Phase’s tier posture using admissible evidence and that durability/recovery trends do not indicate instability that would contaminate Phase 12 verification.
\end{itemize}

\subsection*{12. Risk Controls and Stop Rules}
\phantomsection

\begin{itemize}
\item Stage 0.5 and Stage 1.F remain controlling authority for stop posture and escalation.
\item Phase 11 high-risk exposure controls:
  \begin{itemize}
  \item Maximal \textbf{RUN} and long \textbf{RUCK} are staged to avoid stacking; spacing intent per Stage 3.A is enforced; reslot rather than compress.
  \item Foot integrity is decisive: hotspots and gait drift constrain ruck/run gates and can force deload-only conversion and reslot.
  \item Environment invalidation posture: unsafe footing/heat/air quality triggers postpone or indoor equivalent only when Stage 2 validity/logging preserves comparability.
  \item Underwater governance is absolute: if governance cannot be guaranteed, underwater is not scheduled and cannot earn gate credit.
  \end{itemize}
\end{itemize}

\clearpage
\begin{landscape}

\subsection*{13. Required Capabilities and Dependencies}
\phantomsection

\begingroup
\scriptsize
\setlength{\tabcolsep}{2pt}
\renewcommand{\arraystretch}{1.12}

\begin{longtblr}{
  width=\linewidth,
  rowhead=1,
  colspec = {
    X[l,1.40]
    X[l,2.85]
    Q[l,wd=1.30in]
    Q[l,wd=1.20in]
    X[l,3.25]
  },
  hlines,
  vlines,
  cells = {hsep=6pt, vsep=6pt, halign=l, valign=t},
  cell{1-Z}{1-5} = {fg=black},
  row{odd} = {bg=white},
  row{even} = {bg=LightGray},
  row{1} = {bg=StageBlue!15, fg=black, font=\bfseries, halign=l, valign=t},
}
\Hdr{Dependency \textbf{ID}} &
\Hdr{Required Capability Category and Tier} &
\Hdr{Due-By week-in-phase} &
\Hdr{Owner} &
\Hdr{Fail Action} \\

\textbf{DEP-1C-007} \& \textbf{DEP-1C-008} &
7.11 Footwear Systems and Foot Care — Tier: Phase 11 gate-ready &
Week 1 &
Trainee &
\textbf{HOLD} any \textbf{RUN}/\textbf{RUCK} gate attempt until footwear interface is stable; reslot tests to next protected window; document durability risk flags. \\

\textbf{DEP-1C-013} &
7.07 Load-Bearing Rucking and Carry Systems — Tier: long-ruck gate-ready &
Week 8 &
Equipment Lead &
If not gate-ready, Week 10 long ruck is not attempted; reslot to the next protected window or extend phase per Stage 3.E; no proxy substitution. \\

\textbf{DEP-1C-006} &
7.18 Testing Timing Measurement and Course-Setup Tools — Tier: verified &
Week 1 &
Coach Lead &
If tools/routes/pool verification cannot be stabilized, tests are inadmissible; reslot; do not treat compromised evidence as gate credit. \\

\textbf{DEP-1C-011} &
7.17 Aquatic Training and Water Safety Equipment — Tier: verified (if \textbf{SWIM} is executed) &
Week 5 &
Equipment Lead &
If pool access/verification is unstable, \textbf{SWIM} test is reslotted; do not substitute with non-swim evidence for \textbf{SWIM} gates. \\

\textbf{DEP-1C-015} &
7.09 Health Monitoring Diagnostics and Biometrics — Tier: logging-ready &
Week 1 &
Trainee &
Missing readiness logging triggers conservative downshift; if it prevents admissibility, default \textbf{HOLD} until evidence integrity is restored. \\

\textbf{DEP-1C-017} &
7.12 Recovery Mobility and Tissue-Care Implements — Tier: operational &
Week 1 &
Trainee &
If recovery tools are not available, maintain conservative dose and reslot high-cost tests when durability drift appears; do not force gates under escalating pain. \\

\end{longtblr}

\endgroup

\subsection*{14. Validity Requirements}
\phantomsection

\begin{itemize}
\item Stage 1.F governs admissibility. Gate decisions rely on admissible evidence only.
\item Tests must be tagged \textbf{VALID} or \textbf{INVALID}. \textbf{INVALID} tests provide no gate credit and are treated as not yet tested for gate decisions.
\item \textbf{CAL} tests require compliant artifact evidence for \textbf{VALID} gate use per Stage 2.
\item Environment/tool stability and logging completeness per Stage 2.E are required for comparability; missing required context fields defaults evidence to inadmissible for gate use.
\item Reslot posture per Stage 3.E applies to missed/invalid tests; no stacking make-ups.
\end{itemize}

\end{landscape}
\clearpage

\subsection*{15. Change Control Notes}
\phantomsection

\begin{itemize}
\item Allowed without patch:
  \begin{itemize}
  \item reslotting a missed/invalid test to the next protected window within the phase, preserving Stage 3.A stacking/spacing intent,
  \item swapping order of tests within the same protected week when it preserves validity and does not increase risk.
  \end{itemize}
\item Requires patch:
  \begin{itemize}
  \item adding test weeks, changing phase duration, changing gate test set, changing work:deload pattern.
  \end{itemize}
\item Any change that impacts crosswalk keys or gate logic requires change control per Stage 1.F and Stage 8 impact review posture.
\end{itemize}

\subsection*{16. Compliance Checklist}
\phantomsection

\begin{itemize}
\item Phase \textbf{ID}, name, duration, and governing status present.
\item Weekly structure table includes Cardio cell structure (Modality; Minutes; RPE plus Talk Test; Continuity continuous planned).
\item Cardio minimum and steps requirement stated in body text.
\item Deload/test windows are explicit and align to Stage 3 phase posture.
\item Test battery references Stage 2 instruments only and contains no invented tests.
\item Dependencies use Stage 7 category IDs and owners; fail actions are operational.
\item Validity/admissibility posture stated; invalid evidence provides no gate credit and triggers reslot.
\item Governing Status justification: Draft — Non-Deployable because dependency IDs are placeholders (use Stage 1 dependency \textbf{ID}) in gate-relevant dependency table fields.
\end{itemize}

\end{document}